\documentclass[12pt]{article}
\usepackage[utf8]{inputenc}
\usepackage[T1]{fontenc}
\usepackage[margin=1in]{geometry}
\usepackage{setspace}
\usepackage{xcolor}
\usepackage[colorlinks=true, urlcolor=blue!60!black, linkcolor=blue!60!black]{hyperref}
\setlength{\parindent}{2em}
\setlength{\parskip}{1.25em}
\renewcommand{\baselinestretch}{1.0}
\begin{document}
\begin{center}
{\LARGE \textit{Texas v. Johnson}, 491 U.S. 397 (1989)\par}
\vspace{0.8em}
{\large Case Brief --- Four-Page Edition\par}
\end{center}
\vspace{1em}

\section{Memory Jogger}

Burning an American flag as political protest is expressive conduct protected by the First Amendment, and a State may not criminalize it simply because society finds the message offensive or wishes to preserve the flag as a symbol of national unity (\href{https://www.courtlistener.com/opinion/112304/texas-v-johnson/}{491 U.S. 397}).

\section{Detailed Case Facts}

During the 1984 Republican National Convention in Dallas, Gregory Lee Johnson joined the ``Republican War Chest Tour,'' a march of about 100 demonstrators protesting the Reagan administration and Dallas-based corporations with slogans, literature, and ``die-ins.'' Some protesters spray-painted walls, though Johnson did not; another demonstrator took a flag from a flagpole and handed it to him. At Dallas City Hall, Johnson doused the flag with kerosene and set it ablaze while protesters chanted, ``America, the red, white, and blue, we spit on you.'' No one was injured or threatened, though several onlookers were seriously offended.

Johnson alone was charged, solely with desecrating a venerated object under Texas Penal Code \S~42.09(a)(3), which punished physically mistreating a flag ``in a way that the actor knows will seriously offend one or more persons likely to observe or discover his action.''

\section{Procedural History}

A jury convicted Johnson; he was sentenced to one year in prison and fined \$2,000, and the Dallas Court of Appeals affirmed. The Texas Court of Criminal Appeals reversed, holding the statute unconstitutional as applied: the burning was symbolic speech, the State could not mandate national unity under \textit{Barnette}, and no breach of the peace occurred. The Supreme Court granted certiorari and affirmed on June 21, 1989.

\section{Judicial Votes}

Affirmed, \textbf{5--4}.
\begin{itemize}
\item \textbf{Majority:} Justice Brennan, with Justices Marshall, Blackmun, Scalia, and Kennedy joining.
\item \textbf{Concurrence:} Justice Kennedy, joining the majority in full.
\item \textbf{Dissents:} Chief Justice Rehnquist, joined by Justices White and O'Connor; Justice Stevens.
\end{itemize}

\section{Holding}

The conviction is inconsistent with the First Amendment. Expressive conduct is protected when there is an intent to convey a particularized message and a great likelihood that viewers will understand it (\textit{Spence v. Washington}). When the State regulates such conduct for reasons related to suppressing expression---protecting a symbol's meaning or shielding onlookers from offense---\textit{O'Brien}'s lenient test does not apply, and the restriction faces ``the most exacting scrutiny,'' which a desire to suppress offensive ideas cannot satisfy. The ruling was as-applied; footnote 8 noted that theft, trespass, and arson remain punishable, and the remarks about the fitting response to flag burning are dicta.

\section{Analysis of Opinions}

\textbf{Majority (Brennan).} The burning, at the climax of a march timed to Reagan's renomination, was ``sufficiently imbued with elements of communication,'' as Texas conceded. Texas's breach-of-the-peace interest was not implicated: nothing happened, and offensive speech cannot be presumed to incite violence absent \textit{Brandenburg} incitement or \textit{Chaplinsky} fighting words. Its interest in preserving the flag as a symbol of nationhood arises only when treatment of the flag communicates a message, so it is related to expression. Because the statute punished burning only when it caused ``serious offense,'' while respectful disposal is lawful, exacting scrutiny applied and was not met. Government may not prescribe an orthodox meaning for national symbols, and a flag-only exception would have ``no discernible or defensible boundaries.''

\textbf{Concurrence (Kennedy).} Acknowledging the ``personal toll,'' Kennedy wrote that the Constitution compelled the result, for ``[i]t is poignant but fundamental that the flag protects those who hold it in contempt.''

\textbf{Dissent (Rehnquist).} Recounting the flag's history in poetry, war, and the laws of 48 States, Rehnquist likened flag burning to fighting words, ``the equivalent of an inarticulate grunt or roar,'' noted that Johnson remained free to express his ideas in other ways, and accused the Court of acting as a ``platonic guardian.''

\textbf{Dissent (Stevens).} The flag is an ``important national asset'' symbolizing freedom and tolerance; the statute was content-neutral and imposed only a ``trivial burden'' on speakers with ample alternatives.

\section{Examples: Future Applications}

\begin{enumerate}
\item \textbf{Same-Side (Protected): Flag Protection Act.} A federal ban on burning or defiling any flag, without a ``serious offense'' element, still targets the flag's symbolic meaning and falls, as \textit{United States v. Eichman} held.
\item \textbf{Same-Side (Protected): State Symbols.} A law against burning copies of the state flag or constitution ``in a manner likely to cause public outrage'' targets dissent and triggers strict scrutiny.
\item \textbf{Opposite-Side (Prosecutable): Theft and Open Flame.} A protester who steals a flag, trespasses, and ignites it in a dry park may be prosecuted for theft, trespass, or violating a content-neutral fire ordinance.
\item \textbf{Opposite-Side (Prosecutable): Defacing a Monument.} Spray-painting anti-war slogans on the Lincoln Memorial may be punished under a content-neutral ban protecting public property.
\item \textbf{Fence-Sitter: Confrontation at a Cemetery.} Burning a flag inches from a grieving military family while shouting threats and urging a crowd to attack edges toward fighting words or incitement, so the outcome turns on the record.
\end{enumerate}

\section{Critique}

\textbf{Originalist and textualist views.} Justice Scalia's vote is cited as proof that principle can override personal distaste, but history-and-tradition scholars question whether burning an object falls within the original meaning of ``speech,'' noting Brennan's reliance on modern precedent rather than Founding-era sources.

\textbf{Progressive and doctrinal views.} Geoffrey Stone and John Hart Ely praise \textit{Johnson} for enforcing the line between content-neutral regulation of conduct and content-based restrictions triggered by audience offense. Critics contend that strict content neutrality undervalues civic cohesion by treating a uniquely unifying emblem like any other symbol.

\section{Key Quotations}

\begin{quote}
``If there is a bedrock principle underlying the First Amendment, it is that the government may not prohibit the expression of an idea simply because society finds the idea itself offensive or disagreeable.''
\end{quote}
\noindent --- Justice Brennan, Opinion of the Court, 491 U.S. at 414.

\begin{quote}
``The hard fact is that sometimes we must make decisions we do not like. We make them because they are right, right in the sense that the law and the Constitution, as we see them, compel the result.''
\end{quote}
\noindent --- Justice Kennedy, concurring, 491 U.S. at 420--421.

\begin{quote}
``We do not consecrate the flag by punishing its desecration, for in doing so we dilute the freedom that this cherished emblem represents.''
\end{quote}
\noindent --- Justice Brennan, Opinion of the Court, 491 U.S. at 420.

\end{document}
